\section{Hyperparameters and Configurations}
\label{app:hyper}

Table~\ref{tab:hyper} lists the exact configurations used for both
production runs. Both jobs ran on a single NVIDIA H100 80~GB in bf16 via
the Modal cloud platform, with the full stdout stream teed to a
persistent volume and uploaded to the Hub after (and, on the 7B run,
during) training. The W\&B project \texttt{knowing\_using\_gap} holds
per-step and per-epoch curves, LoRA/probe matrix panels (every 5
epochs), and locked checkpoint artifacts.

\begin{table}[h]
\centering
\small
\begin{tabular*}{\columnwidth}{@{\extracolsep{\fill}} l l l}
\toprule
Parameter & 0.5B run & 7B run \\
\midrule
Model & Qwen2.5-0.5B-Instruct & Qwen2.5-7B-Instruct \\
LoRA $r$ / $\alpha$ / dropout & 16 / 32 / 0.05 & 16 / 32 / 0.05 \\
LoRA targets & q,k,v,o,gate,up,down & q,k,v,o,gate,up,down \\
Probe layer $l^\star$ & 12 of 24 & 14 of 28 \\
Probe head & $896 \to 1600$ & $3584 \to 1600$ \\
Backbone LR & $2{\times}10^{-5}$ & $2{\times}10^{-5}$ \\
Probe LR & $10^{-3}$ & $10^{-3}$ \\
Schedule & cosine, warmup $10\%$ & cosine, warmup $10\%$ \\
Batch (dev $\times$ accum) & $4 \times 4$ & $4 \times 4$ \\
Steps per epoch & 75 & 75 \\
Epochs (planned) & 50 & 50 \\
Epochs completed & 50 & 27 \\
Gate metric / $\tau_0$ & max / 0.8 & max / 0.8 \\
$\lambda_0$ / $\delta_0$ & 0.3 / 0.1 & 0.3 / 0.1 \\
Controller decay $\gamma$ & 0.7 & 0.7 \\
Damping floor $m_{\min}$ & 0.5 & 0.5 \\
Boost factor / cap & $1.5$ / $8\times$ & $1.5$ / $8\times$ \\
Lifecycle pushes & every 10 ep & every 10 ep \\
Matrices to W\&B & every 5 ep & every 5 ep \\
\bottomrule
\end{tabular*}
\caption{Full hyperparameter configurations.}
\label{tab:hyper}
\end{table}

\section{Dataset Construction}
\label{app:dataset}

The generator (released as
\texttt{scripts/generate\_synthetic\_dataset.py}) proceeds in three
steps. \textbf{(1) Entity vocabulary:} 1{,}600 invented entity names are
generated from syllable templates; each is assigned an id in
$[0, 1600)$ for the probe target space. \textbf{(2) Memorization split
$\mathcal{D}_{\mathrm{mem}}$:} 1{,}200 question--answer pairs in a fixed
interrogative template ending with \texttt{Answer:} \emph{label}; the
suffix is present during CE training (so the format is learned) and
stripped during all evaluations (so it cannot leak the answer into the
prompt). \textbf{(3) Generalization split $\mathcal{D}_{\mathrm{gen}}$}
(450 prompts): 200 \emph{chaining} prompts combine two trained facts
into a compositional question, and 250 \emph{intersection} prompts
re-ask trained facts with paraphrased templates and surface forms never
seen in training. Free-generation evaluations use left-padding,
greedy decoding, and strict exact match after answer extraction; the
likelihood evaluation \eqref{eq:lf} scores the first answer token and
the token-level overlap under teacher forcing. The per-epoch gate
evaluations run on a fixed random subsample of 208 memorization
prompts (13 batches) plus the full $\mathcal{D}_{\mathrm{gen}}$, keeping
epoch overhead at $\sim$1--2 minutes.

\section{Full Per-Epoch Curves}
\label{app:curves}

Table~\ref{tab:curves05} reports the 0.5B run every 5 epochs and
Table~\ref{tab:curves7} the 7B run every 3 epochs (the complete
per-epoch JSON histories are in the released telemetry).
Eval loss is the held-out generalization-split loss logged by the
trainer each epoch; $m(t)$ is the controller's cumulative LR
multiplier.

\begin{table}[h]
\centering
\scriptsize
\setlength{\tabcolsep}{3.2pt}
\begin{tabular*}{\columnwidth}{@{\extracolsep{\fill}} r rrrrrr}
\toprule
ep & $\amemlf$ & tok.\ & $\amemem$ & $\agen$ & $\lambda$ & eval \\
\midrule
1 & 0.000 & 0.337 & 0.000 & 0.000 & 0.000 & 5.03 \\
5 & 0.000 & 0.748 & 0.000 & 0.000 & 0.000 & 3.96 \\
10 & 0.010 & 0.758 & 0.000 & 0.000 & 0.000 & 4.23 \\
15 & 0.025 & 0.771 & 0.000 & 0.000 & 0.000 & 4.22 \\
20 & 0.030 & 0.786 & 0.005 & 0.000 & 0.000 & 4.40 \\
25 & 0.125 & 0.819 & 0.005 & 0.000 & 0.000 & 4.45 \\
30 & 0.240 & 0.857 & 0.055 & 0.000 & 0.000 & 4.50 \\
35 & 0.275 & 0.866 & 0.010 & 0.000 & 0.000 & 4.55 \\
40 & 0.355 & 0.895 & 0.050 & 0.000 & 0.000 & 4.57 \\
45 & 0.375 & 0.899 & 0.010 & 0.000 & 0.000 & 4.61 \\
50 & 0.380 & 0.900 & 0.025 & 0.000 & 0.000 & 4.64 \\
\bottomrule
\end{tabular*}
\caption{0.5B run (complete, 3{,}750 steps). $m(t)$: $1.0$ (ep 1--5),
$0.7$ (ep 6--7), $0.5$ (ep 8--50); one memorization diagnosis at the
controller's stagnation trigger. The gate never opens.}
\label{tab:curves05}
\end{table}

\begin{table}[h]
\centering
\scriptsize
\setlength{\tabcolsep}{3.2pt}
\begin{tabular*}{\columnwidth}{@{\extracolsep{\fill}} r rrrrrr}
\toprule
ep & $\amemlf$ & tok.\ & $\amemem$ & $\agen$ & $\lambda$ & eval \\
\midrule
1 & 0.000 & 0.386 & 0.000 & 0.000 & 0.000 & 5.10 \\
3 & 0.015 & 0.752 & 0.000 & 0.000 & 0.000 & 3.21 \\
6 & 0.005 & 0.760 & 0.000 & 0.000 & 0.000 & 4.02 \\
9 & 0.025 & 0.783 & 0.000 & 0.000 & 0.000 & 4.06 \\
12 & 0.110 & 0.822 & 0.005 & 0.000 & 0.000 & 4.19 \\
15 & 0.320 & 0.895 & 0.015 & 0.000 & 0.000 & 4.24 \\
18 & 0.570 & 0.946 & 0.030 & 0.000 & 0.000 & 4.34 \\
21 & 0.745 & 0.971 & 0.030 & 0.000 & 0.000 & 4.40 \\
24 & \textbf{0.805} & 0.979 & 0.055 & 0.000 & 0.015 & 4.44 \\
25 & 0.825 & 0.981 & 0.025 & 0.000 & 0.075 & 4.41 \\
27 & 0.875 & 0.987 & 0.020 & 0.000 & 0.225 & --- \\
\bottomrule
\end{tabular*}
\caption{7B run (27 completed epochs; terminated during epoch 28 by
compute-budget exhaustion, hence no epoch-28 eval). The crossing of
$\tau_0{=}0.8$ at epoch 24 opens the gate; $\lambda(t)$ ramps to
$0.225$ by epoch 27. Epoch-27 values were read from the live run log.}
\label{tab:curves7}
\end{table}

\section{The Decay-Spiral Failure and Its Repair}
\label{app:spiral}

\textbf{Failure.} In a pilot 50-epoch 0.5B run (v3), the eval-aware
controller issued eval-driven decays every two epochs for five
consecutive episodes while held-out loss diverged, driving the
cumulative LR multiplier along
$1.0 \to 0.49 \to 0.34 \to 0.24 \to 0.17$, breaching the controller's
own nominal minimum ($0.3$) because the eval-driven pathway applied
decays without consulting the instability path's floor. Memorization
was choked: likelihood recall stayed pinned near $0.02$ and the gate
threshold became unreachable. The run was stopped to avoid burning
compute on dynamics that the controller itself had frozen.

\textbf{Diagnosis.} The failure is structural, not incidental: in the
memorization regime, held-out divergence is the \emph{expected}
signature that the saturation gate waits for, so an eval-driven
damping rule without a floor is guaranteed to eventually fight the
gate.

\textbf{Repair.} The eval-driven pathway now clamps the cumulative
multiplier to $m_{\min}{=}0.5$ (Eq.~\ref{eq:floor}) and refuses
further eval-driven decays at the floor. Validation proceeded in
three stages: (i) unit tests on a stubbed trainer assert the floor is
respected exactly ($m$ lands on $0.5$, never below; optimizer LRs
match base $\times$ floor) among 138 checks; (ii) a dedicated 5-epoch
GPU run with an unreachable gate reproduced the failing regime
($\lambda{=}0$ throughout, eval rising) and showed one eval-driven
decay ($1.0 \to 0.7$) with the floor never breached; (iii) both
production runs in this paper clamp at $0.5$ and hold
(Fig.~\ref{fig:dyn}c) while memorization proceeds unimpeded to
saturation at 7B.

\section{Artifact Map and Log Reconstruction}
\label{app:artifacts}

\textbf{Hugging Face Hub (user \texttt{codewithdark}).}
\texttt{mlpr-qwen2.5-0.5b-instruct-50ep-adaptive-v4}: final adapter,
probe head (\texttt{probe\_head.pt}), tokenizer, and full telemetry
(\texttt{gate\_history.json}, \texttt{adaptive\_history.json},
\texttt{trainer\_state.json} with 1{,}925 log entries,
\texttt{training\_log.txt}, 1.5\,MB). The 7B repository
\texttt{mlpr-qwen2.5-7b-instruct-50ep-adaptive-v4} holds the epoch-26
adapter (the trainer's last in-progress hub commit) and the training
log through epoch 26; its per-epoch gate and controller records for
epochs 1--25 are parsed from the log's structured
\texttt{[MEM\ GATE]} / \texttt{[ADAPTIVE]} lines, and the epoch-27
record was captured verbatim from the live log stream before
termination. \textbf{Weights \& Biases} (project
\texttt{knowing\_using\_gap}): runs
\texttt{mlpr-qwen2.5-0.5b-50ep-adaptive-v4} and
\texttt{mlpr-qwen2.5-7b-50ep-adaptive-v4} with per-step losses,
learning rates, gradient norms, per-epoch adaptive decisions, LoRA/probe
matrix panels every 5 epochs (168 heatmaps each of $B$, $A$, and
$\Delta W$), and locked checkpoint artifacts.
\textbf{Codebase}: the released zip contains the full training stack
(\texttt{src/} with probe, LoRA setup, gate callback, adaptive
controller, likelihood evaluator, W\&B matrix logger), the Modal
deployment (\texttt{modal\_train.py}), the dataset generator, the
analysis and figure scripts, and this paper under \texttt{paper/}.

\section{The Likelihood Recall Metric in Detail}
\label{app:lfmetric}

For a memorization pair $(x, y)$ the prompt is built by stripping the
\texttt{Answer:} suffix from the training text. The metric performs a
single teacher-forced forward pass over
$\texttt{prompt} \oplus \texttt{answer}$ with left padding and records
(i) strict first-answer-token exact match \eqref{eq:lf}, and (ii) the
fraction of answer tokens whose greedy next-token prediction matches
the reference (token recall). Both are macro-averaged over the
subsample. Compared with free-generation EM, this removes decoding
length, format, and repetition effects; the residual gap between token
recall and strict EM isolates single-token identity from multi-token
assembly. Substring containment of the target in the free generation
complements both by detecting \emph{stored but misframed} answers---the
$76.5\%$ vs.\ $2\%$ contrast in Table~\ref{tab:main}.
